\documentclass[12pt]{article}

\usepackage[margin=1in]{geometry}
\usepackage[T1]{fontenc}
\usepackage{lmodern}
\usepackage{microtype}
\usepackage{amsmath,amssymb,mathtools}
\usepackage{booktabs}
\usepackage{setspace}
\usepackage{hyperref}
\usepackage{xurl}

\hypersetup{
  colorlinks=true,
  linkcolor=black,
  citecolor=black,
  urlcolor=blue
}

\onehalfspacing
\widowpenalty=10000
\clubpenalty=10000

\title{The Final Parsec Equilibrium:\\
A Rate Crossover Between Stellar Hardening and Gravitational Waves}
\author{Light-Meter Physics}
\date{}

\begin{document}

\maketitle

\begin{abstract}
A massive black-hole binary can shrink through stellar encounters at large separations and gravitational-wave emission at sufficiently small separations. This paper constructs a transparent circular-binary model in which the stellar hardening rate is held constant and the leading-order quadrupole gravitational-wave rate is added. The two inward rates are equal at a characteristic separation $a_{\mathrm{eq}}=(B/s)^{1/5}$, where $s$ is the stellar hardening parameter and $B$ sets the gravitational-wave decay rate. This ``equilibrium'' is a rate crossover, not a stationary orbit or a stall: both channels continue to shrink the binary. A worked equal-mass example with $10^8M_\odot$ components and an explicitly maintained stellar background gives $a_{\mathrm{eq}}\simeq0.0192\ \mathrm{pc}$, a separation light-time of about 22.9 days, and equal channel-by-channel local shrinking times of about 158 Myr. These local times are not a predicted stalling duration or, by themselves, the total merger time. The paper distinguishes instantaneous shrinking times from integrated gravitational-wave-only merger times, explains why stellar-supply depletion makes stalling model-dependent, and states the limitations of the fixed-background, circular, Newtonian/quadrupole approximation. The crossover is useful bookkeeping within its assumptions; it is not a universal solution to the final-parsec problem.
\end{abstract}

\noindent\textbf{Keywords:} massive black-hole binaries, final-parsec problem, stellar hardening, loss cone, gravitational waves, rate crossover

\section{Introduction}

When two massive black holes become a bound binary after a galaxy merger, the binary must lose orbital energy and angular momentum to reach coalescence. At relatively large separations, interactions with the stellar background can harden the binary. At sufficiently small separations, gravitational-wave emission drives rapid orbital decay. Whether stellar encounters continue to supply enough orbital energy loss as the binary tightens is central to the so-called final-parsec problem \cite{Quinlan,Gualandris,Vasiliev}.

The purpose of this paper is narrower than predicting the complete evolution of a galactic nucleus. We compare two idealized contributions to the rate of change of a circular binary's separation: a stellar hardening law with constant hardening parameter and the leading-order quadrupole gravitational-wave law. Their rates cross at a well-defined separation. We call that scale the equal-rate separation, $a_{\mathrm{eq}}$.

The word ``equilibrium'' in the title refers only to equality of the two inward shrinking rates at that separation. The orbit is not stationary there: both rates have the same sign, so the combined model continues to shrink the binary. Nor does the crossover prescribe how long the binary may wait for stellar encounters before reaching it. The answer to that question depends on the evolving stellar supply and other environmental physics.

We define the orbital model, derive each rate and their crossover, and evaluate a numerical example. We then distinguish local from integrated timescales and summarize the model limits.

\section{Binary variables and orbital energy}

For fixed component masses $m_1$ and $m_2$, define the total mass, reduced mass, and symmetric mass ratio by
\begin{equation}
  M_{\mathrm{tot}}\equiv m_1+m_2,
  \qquad
  \mu\equiv\frac{m_1m_2}{M_{\mathrm{tot}}},
  \qquad
  \eta\equiv\frac{m_1m_2}{M_{\mathrm{tot}}^2}.
  \label{eq:mass-variables}
\end{equation}
Let $a$ denote the relative separation of the components. For a Newtonian circular relative orbit,
\begin{equation}
  E_{\mathrm{orb}}=-\frac{Gm_1m_2}{2a},
  \qquad
  L_{\mathrm{orb}}=\mu\sqrt{GM_{\mathrm{tot}}a}.
  \label{eq:orbital-energy-angular-momentum}
\end{equation}
Here $M_{\mathrm{tot}}$ is the binary's total mass and $L_{\mathrm{orb}}$ its orbital angular momentum. In this paper $a$ denotes binary separation, not the Kerr spin length parameter used in Paper 19. In an equal-mass binary, each component is a distance $a/2$ from the barycenter.

The time $t$ is the source-frame secular evolution parameter of this Newtonian binary model, and overdots denote derivatives with respect to $t$. It is not the proper time of either black hole. We hold the component masses fixed, neglect mass growth, and impose $e=0$ throughout the orbit-averaged model. This prescribes a circular sequence; the scalar stellar-hardening law does not determine eccentricity evolution or demonstrate that stellar encounters maintain circularity. Eccentric binaries require additional evolution equations and generally have different gravitational-wave decay rates.

\section{Stellar hardening with a maintained supply}

For a sufficiently hard binary embedded in a stellar background, the dimensionless hardening coefficient $H$ is conventionally used to write the rate of increase of inverse separation as \cite{Quinlan}
\begin{equation}
  s\equiv\left.\frac{\mathrm{d}(1/a)}{\mathrm{d}t}\right|_{\star}
  =-\frac{\dot a_\star}{a^2}
  \simeq H\frac{G\rho}{\sigma_\star},
  \label{eq:hardening-rate}
\end{equation}
Here $\rho$ and $\sigma_\star$ are the prescribed stellar-reservoir mass density and one-dimensional velocity dispersion entering the scattering model. They are not calculated self-consistently as the reservoir is depleted. The value $H=15$ used below is illustrative, not a universal exact coefficient. The units of $s$ are inverse length per unit time. When $H$, $\rho$, and $\sigma_\star$ are treated as constant, $s$ is constant as well.

The chain rule for the inverse separation is
\begin{equation}
  \frac{\mathrm{d}(1/a)}{\mathrm{d}t}
  =-\frac{\dot a}{a^2}.
\end{equation}
Applying it to the stellar contribution gives
\begin{equation}
  \dot a_\star=-sa^2.
  \label{eq:stellar-shrink-rate}
\end{equation}
The negative sign denotes shrinking separation. The corresponding instantaneous local shrinking time is
\begin{equation}
  t_\star(a)\equiv\frac{a}{|\dot a_\star|}
  =\frac{1}{sa}.
  \label{eq:stellar-local-time}
\end{equation}
This time depends on the current separation; it is not a fixed duration for the entire stellar-hardening phase.

The constant-$s$ approximation is a maintained-supply model. It presumes that stars continue to reach orbits passing close enough to interact with the binary. In a real nucleus, the density, velocity distribution, binary eccentricity, and supply of such stars can evolve. Equation~\eqref{eq:hardening-rate} is therefore a useful local model under specified assumptions, not a universal law for every galactic environment.

\section{Stellar supply and conditional stalling}

The loss cone is the region of stellar orbital phase space occupied by orbits that pass sufficiently close to the binary for strong interaction. If those stars are ejected faster than other stars can enter the loss cone, the supply of new encounters can diminish and the stellar hardening rate can fall. This is the origin of the idealized stalling concern in spherical stellar systems. Studies of non-spherical systems and collisionless loss-cone refilling show why the outcome depends on galactic structure and stellar dynamics rather than on a single universal waiting time \cite{Gualandris,Vasiliev}.

The fixed-background model in Eq.~\eqref{eq:hardening-rate} does not calculate depletion, refilling, or a stall duration. It also does not imply that every binary stalls: the supply can be replenished in some systems, while other systems may have a substantially reduced stellar hardening rate. Any estimate of the time spent in a low-supply phase would require a model for the galaxy and its evolving stellar orbits.

\section{Gravitational-wave decay for a circular binary}

In the leading-order quadrupole approximation, the separation of an isolated circular binary evolves according to \cite{Peters}
\begin{equation}
  \dot a_{\mathrm{GW}}=-\frac{B}{a^3},
  \qquad
  B\equiv\frac{64}{5}
  \frac{G^3m_1m_2M_{\mathrm{tot}}}{c^5}.
  \label{eq:gw-decay}
\end{equation}
The coefficient $B$ has dimensions of length to the fourth power per time. The steep $a^{-3}$ dependence means that halving the separation increases the magnitude of the instantaneous gravitational-wave shrinking rate by a factor of eight.

Two timescales follow from this rate but answer different questions. The instantaneous local shrinking time is
\begin{equation}
  t_{a,\mathrm{GW}}(a)\equiv\frac{a}{|\dot a_{\mathrm{GW}}|}
  =\frac{a^4}{B}.
  \label{eq:gw-local-time}
\end{equation}
If the binary then evolves under gravitational waves alone, integrating Eq.~\eqref{eq:gw-decay} from an initial separation $a$ to zero gives
\begin{equation}
  t_{\mathrm{merge,GW}}(a)
  =\int_0^a\frac{a'^3}{B}\,\mathrm{d}a'
  =\frac{a^4}{4B}.
  \label{eq:gw-merger-time}
\end{equation}
For a finite endpoint $a_f<a_i$, the same approximation gives
\begin{equation}
  \Delta t_{\mathrm{GW}}(a_i\to a_f)
  =\frac{a_i^4-a_f^4}{4B}.
  \label{eq:gw-finite-endpoint}
\end{equation}
Setting $a_f=0$ yields a formal circular-quadrupole coalescence estimate; the approximation does not describe the final strong-field plunge and merger. The factor of four distinguishes this integrated estimate from the instantaneous shrinking time. Neither quantity includes continued stellar hardening.

\section{The equal-rate separation}

The magnitudes of the stellar and gravitational-wave contributions are equal when
\begin{equation}
  |\dot a_\star|=|\dot a_{\mathrm{GW}}|,
  \qquad
  sa_{\mathrm{eq}}^2=\frac{B}{a_{\mathrm{eq}}^3}.
  \label{eq:equal-rates}
\end{equation}
Solving gives the equal-rate separation
\begin{equation}
  \boxed{a_{\mathrm{eq}}=\left(\frac{B}{s}\right)^{1/5}.}
  \label{eq:equal-separation}
\end{equation}
The relative importance of the two terms at any separation can be written
\begin{equation}
  \frac{|\dot a_{\mathrm{GW}}|}{|\dot a_\star|}
  =\frac{B}{sa^5}
  =\left(\frac{a_{\mathrm{eq}}}{a}\right)^5.
  \label{eq:rate-ratio}
\end{equation}
Thus, within the constant-$s$ model, stellar hardening is faster for $a>a_{\mathrm{eq}}$, the two rates match at $a=a_{\mathrm{eq}}$, and gravitational-wave decay is faster for $a<a_{\mathrm{eq}}$.

Crossover values in the literature can differ when the comparison criterion differs. Equating the stellar instantaneous shrinking time $1/(sa)$ to the GW-only integrated time $a^4/(4B)$ gives $a_{\mathrm Q}=(4B/s)^{1/5}=4^{1/5}a_{\mathrm{eq}}$. This is a different criterion from the equal-instantaneous-rate definition used here; it does not change Eq.~\eqref{eq:equal-separation} \cite{Quinlan}.

At every positive separation, the combined rate is
\begin{equation}
  \dot a=\dot a_\star+\dot a_{\mathrm{GW}}
  =-sa^2-\frac{B}{a^3}<0.
  \label{eq:combined-rate}
\end{equation}
The total inverse-separation hardening rate is therefore
\begin{equation}
  \frac{\mathrm{d}(1/a)}{\mathrm{d}t}
  =-\frac{\dot a}{a^2}
  =s+\frac{B}{a^5}.
  \label{eq:total-hardening-rate}
\end{equation}
At $a=a_{\mathrm{eq}}$ it equals $2s$. Both channels shrink the orbit; neither reverses the inspiral at the crossover. In this limited sense, ``equilibrium'' means only equal magnitudes of two rates, not force balance, a stable separation, or a static mechanical equilibrium.

At the crossover, the two channel-by-channel local times are equal:
\begin{equation}
  t_\star(a_{\mathrm{eq}})
  =t_{a,\mathrm{GW}}(a_{\mathrm{eq}})
  =\frac{1}{s a_{\mathrm{eq}}}
  \equiv t_0.
  \label{eq:equal-local-times}
\end{equation}
The local time based on the combined instantaneous rate is half this value,
\begin{equation}
  t_{a,\mathrm{total}}(a_{\mathrm{eq}})
  \equiv\frac{a_{\mathrm{eq}}}{|\dot a(a_{\mathrm{eq}})|}
  =\frac{1}{2s a_{\mathrm{eq}}}.
  \label{eq:combined-local-time}
\end{equation}
This combined instantaneous time is not the integrated duration from the crossover onward. To obtain an elapsed time under both channels, define $x\equiv a/a_{\mathrm{eq}}$ and $t_0\equiv1/(s a_{\mathrm{eq}})$. Since $B=s a_{\mathrm{eq}}^5$, Eq.~\eqref{eq:combined-rate} becomes
\begin{equation}
  \frac{\mathrm{d}x}{\mathrm{d}(t/t_0)}
  =-(x^2+x^{-3}).
\end{equation}
Integration from $a_i$ to a smaller endpoint $a_f$ gives
\begin{equation}
  \Delta t(a_i\to a_f)=t_0
  \int_{a_f/a_{\mathrm{eq}}}^{a_i/a_{\mathrm{eq}}}
  \frac{x^3}{1+x^5}\,\mathrm{d}x.
  \label{eq:combined-elapsed-time}
\end{equation}
This assumes the constant stellar contribution continues over the full interval. Taking $a_f=0$ is only a formal endpoint of the Newtonian/quadrupole model, not a calculation of the final strong-field merger.

\section{A worked equal-mass example}

For illustration, take equal black-hole masses and the prescribed stellar-reservoir parameters in Table~\ref{tab:inputs}. These are model inputs, not measurements of a named galaxy. The adopted constants and conversions make the calculation reproducible; the result is not a universal forecast for massive black-hole binaries.

\begin{table}[ht]
  \centering
  \caption{Inputs for the illustrative circular-binary calculation.}
  \label{tab:inputs}
  \begin{tabular}{lc}
    \toprule
    Quantity & Adopted value \\
    \midrule
    Component masses, $m_1=m_2$ & $10^8 M_\odot$ \\
    Stellar mass density, $\rho$ & $10^3 M_\odot\,\mathrm{pc}^{-3}$ \\
    One-dimensional velocity dispersion, $\sigma_\star$ & $200\ \mathrm{km\,s^{-1}}$ \\
    Hardening coefficient, $H$ & $15$ \\
    Gravitational constant, $G$ & $6.67430\times10^{-11}\ \mathrm{m^3\,kg^{-1}\,s^{-2}}$ \\
    Speed of light, $c$ & $299\,792\,458\ \mathrm{m\,s^{-1}}$ \\
    Solar mass, $M_\odot$ & $1.98847\times10^{30}\ \mathrm{kg}$ \\
    Parsec & $3.085677581491367\times10^{16}\ \mathrm{m}$ \\
    Julian year & $31\,557\,600\ \mathrm{s}$ \\
    \bottomrule
  \end{tabular}
\end{table}

Substituting the stellar parameters into Eq.~\eqref{eq:hardening-rate} gives
\begin{equation}
  s\simeq0.3299\ \mathrm{pc^{-1}\,Myr^{-1}}.
  \label{eq:example-hardening}
\end{equation}
With $m_1=m_2=10^8M_\odot$, Eq.~\eqref{eq:equal-separation} then gives
\begin{equation}
  a_{\mathrm{eq}}\simeq0.01920\ \mathrm{pc}.
  \label{eq:example-separation}
\end{equation}
At crossover, $t_0$ is the instantaneous shrinking time for either channel separately; the combined instantaneous time is $t_0/2$. The GW-only formal remaining time from crossover is $t_0/4$. For both channels maintained from crossover to the formal zero-separation endpoint, Eq.~\eqref{eq:combined-elapsed-time} gives
\begin{equation}
  t_{\mathrm{rem,both}}
  =t_0\int_0^1\frac{x^3}{1+x^5}\,\mathrm{d}x
  \simeq0.18065t_0.
  \label{eq:combined-remaining-time}
\end{equation}
The values below distinguish instantaneous rates from formal integrated remaining times; none is a stellar-supply waiting time or an observed merger lifetime.

The same circular orbit has source-frame dominant quadrupole gravitational-wave frequency
\begin{equation}
  f_{\mathrm{GW}}=2f_{\mathrm{orb}}
  =\frac{1}{\pi}\sqrt{\frac{GM_{\mathrm{tot}}}{a^3}},
  \label{eq:gw-frequency}
\end{equation}
where $f_{\mathrm{orb}}=(2\pi)^{-1}\sqrt{GM_{\mathrm{tot}}/a^3}$ is the Newtonian orbital frequency. At $a_{\mathrm{eq}}$, the example gives $f_{\mathrm{GW}}\simeq3.60\ \mathrm{nHz}$.

\begin{table}[ht]
  \centering
  \caption{Instantaneous and formally integrated times at the equal-rate separation.}
  \label{tab:crossover-times}
  \begin{tabular}{lll}
    \toprule
    Quantity & Type and assumptions & Value \\
    \midrule
    Either channel alone & Instantaneous & $157.9\ \mathrm{Myr}$ \\
    Both channels & Instantaneous, summed rate & $78.95\ \mathrm{Myr}$ \\
    GW only & Integrated to formal $a_f=0$ & $39.47\ \mathrm{Myr}$ \\
    Stellar plus GW & Integrated to formal $a_f=0$; constant $s$ & $28.52\ \mathrm{Myr}$ \\
    \bottomrule
  \end{tabular}
\end{table}
The final two entries extrapolate the Newtonian/quadrupole equations to zero separation; they are formal model integrals, not strong-field merger times. The combined entry additionally assumes the maintained stellar supply persists below crossover.

For contrast, the isolated circular GW-only formal remaining time from $a=1\ \mathrm{pc}$ is about $2.91\times10^{14}\ \mathrm{yr}$ in this example. The formal zero-separation endpoint does not model the final plunge and merger. This number illustrates that environmental hardening is important at large separations; it is not the predicted lifetime of the binary in a galaxy, where stellar interactions and the environment evolve.

\section{The crossover in light-time variables}

The gravitational timescale and symmetric mass ratio are
\begin{equation}
  \tau_g\equiv\frac{GM_{\mathrm{tot}}}{c^3},
  \qquad
  \eta\equiv\frac{m_1m_2}{M_{\mathrm{tot}}^2}.
  \label{eq:time-variables}
\end{equation}
Since $m_1m_2=\eta M_{\mathrm{tot}}^2$, the coefficient in Eq.~\eqref{eq:gw-decay} can be rewritten as
\begin{equation}
  B=\frac{64}{5}\eta c^4\tau_g^3.
  \label{eq:B-time-form}
\end{equation}
Define the separation light-time $\tau_L\equiv a/c$. Substituting $a_{\mathrm{eq}}=c\tau_{L,\mathrm{eq}}$ into the equal-rate condition yields
\begin{equation}
  \boxed{\tau_{L,\mathrm{eq}}^5
  =\frac{64}{5}\frac{\eta\tau_g^3}{cs}.}
  \label{eq:lighttime-crossover}
\end{equation}
For the equal-mass example, $\eta=1/4$, $\tau_g\simeq985\ \mathrm{s}$, and
\begin{equation}
  \tau_{L,\mathrm{eq}}\simeq22.9\ \mathrm{days}.
  \label{eq:example-lighttime}
\end{equation}
This is the light-crossing-time equivalent of the separation. It is a convenient way to express a length as a time; it is not the time the binary spends transitioning between mechanisms, the binary's merger time, or a measured signal delay. The scale $\tau_g$ uses the total binary mass as a convenient parameter and does not treat the binary as a single black hole.

\section{Interpretation and model limits}

The equal-rate separation is determined jointly by the binary masses and the adopted stellar hardening rate. Increasing $s$ moves the crossover to a smaller separation, while increasing $B$ moves it outward. This follows directly from $a_{\mathrm{eq}}=(B/s)^{1/5}$. The fifth root makes the crossover comparatively insensitive to moderate changes in either parameter, but it does not remove the dependence on the model inputs.

Near crossover, the conventional hard-binary scale for these inputs is
\begin{equation}
  a_h\equiv\frac{G\mu}{4\sigma_\star^2}
  \simeq1.34\ \mathrm{pc},
  \qquad
  \frac{a_{\mathrm{eq}}}{a_h}\simeq0.0143.
  \label{eq:hard-binary-scale}
\end{equation}
The orbit is also weak-field and nonrelativistic at crossover:
\begin{equation}
  \frac{GM_{\mathrm{tot}}}{a_{\mathrm{eq}}c^2}\simeq4.99\times10^{-4},
  \qquad
  \frac{v_{\mathrm{rel}}}{c}
  =\sqrt{\frac{GM_{\mathrm{tot}}}{a_{\mathrm{eq}}c^2}}
  \simeq0.0223.
  \label{eq:crossover-regime}
\end{equation}
The circular orbital period is about $17.6\ \mathrm{yr}$, far shorter than the channel-by-channel shrinking time. These checks support use of a hard-binary, weak-field, orbit-averaged description near crossover. They do not validate a maintained stellar supply, and the weak-field quadrupole approximation eventually fails in the strong-field regime.

The calculation has a specific domain:
\begin{itemize}
  \item The binary is treated as circular, with Newtonian orbital energy and frequency.
  \item Gravitational-wave decay uses the leading-order quadrupole result for an isolated circular binary.
  \item Stellar hardening uses a constant $H$, $\rho$, and $\sigma_\star$, representing a maintained stellar supply.
  \item The worked values are illustrative and do not model a particular galaxy or a self-consistent loss-cone history.
  \item The GW-only merger times exclude stellar hardening by construction; a full history requires evolving the coupled orbit and environment.
\end{itemize}

Eccentricity can change gravitational-wave evolution substantially, and changing stellar structure can make $s$ vary with time and separation. Gas, triaxiality, additional black holes, and other environmental effects can alter the evolution as well. These possibilities are reasons to extend the model when needed, not corrections that can be inferred from the two-rate algebra alone.

\section{Conclusion}

Combining a constant stellar hardening law, $\dot a_\star=-sa^2$, with the leading circular quadrupole law, $\dot a_{\mathrm{GW}}=-B/a^3$, defines an equal-rate separation $a_{\mathrm{eq}}=(B/s)^{1/5}$. At that point each channel contributes equally to a still-negative total $\dot a$. The term ``equilibrium'' therefore denotes a rate crossover, not a stationary or stable orbit and not the onset or duration of a stall.

For the stated equal-mass example, the crossover is near $0.0192\ \mathrm{pc}$, with a separation light-time of about 22.9 days. Each channel's local shrinking time there is about 158 Myr, whereas their combined local shrinking time is about 79 Myr and the isolated GW-only integrated time from that separation is about 39 Myr. Their distinctions are physically and mathematically important. The constant stellar-supply assumption makes the crossover calculable, but the binary's route to it remains dependent on the stellar environment. The model is thus a compact rate comparison, not a universal resolution of the final-parsec problem.

\begin{thebibliography}{9}
\small
\setstretch{1.0}
\raggedright
\urlstyle{same}

\bibitem{Quinlan}
G. D. Quinlan,
``The dynamical evolution of massive black hole binaries---I. Hardening in a fixed stellar background,''
\emph{New Astronomy} \textbf{1}, 35--56 (1996),
\href{https://doi.org/10.1016/S1384-1076(96)00003-6}{doi:\nolinkurl{10.1016/S1384-1076(96)00003-6}}.

\bibitem{Peters}
P. C. Peters,
``Gravitational radiation and the motion of two point masses,''
\emph{Physical Review} \textbf{136}, B1224--B1232 (1964),
\href{https://doi.org/10.1103/PhysRev.136.B1224}{doi:\nolinkurl{10.1103/PhysRev.136.B1224}}.

\bibitem{Gualandris}
A. Gualandris, J. I. Read, W. Dehnen, and E. Bortolas,
``Collisionless loss-cone refilling: there is no final parsec problem,''
\emph{Monthly Notices of the Royal Astronomical Society} \textbf{464}, 2301--2310 (2017),
\href{https://doi.org/10.1093/mnras/stw2528}{doi:\nolinkurl{10.1093/mnras/stw2528}}.

\bibitem{Vasiliev}
E. Vasiliev, F. Antonini, and D. Merritt,
``The final-parsec problem in the collisionless limit,''
\emph{The Astrophysical Journal} \textbf{810}, 49 (2015),
\href{https://doi.org/10.1088/0004-637X/810/1/49}{doi:\nolinkurl{10.1088/0004-637X/810/1/49}}.

\end{thebibliography}

\end{document}